% GENERUOJAMA: python -m src.eksperimentai.jautrumas
% Ranka NELIESTI.
\begingroup
\footnotesize
\setlength{\tabcolsep}{4pt}
\begin{tabularx}{\textwidth}{@{}>{\raggedright\arraybackslash}p{5.4cm}>{\centering\arraybackslash}p{1.6cm}>{\raggedright\arraybackslash}X@{}}
\toprule
\textbf{Palyginimas} & \textbf{Skirtumas} & \textbf{Kada apsiverstų} \\
\midrule
XGBoost prieš Random Forest & $+$1,15 & \textbf{Niekada.} Pirmasis ne blogesnis pagal visus keturis kriterijus, todėl jokie svoriai rezultato nekeičia \\
\addlinespace[2pt]
Random Forest prieš MLP & $+$0,30 & \textbf{Niekada.} Pirmasis ne blogesnis pagal visus keturis kriterijus, todėl jokie svoriai rezultato nekeičia \\
\addlinespace[2pt]
XGBoost prieš Sprendimų medis & $+$0,90 & interpretuojamumo svoris turėtų pasiekti 0,42 (dabar 0,15) \\
\addlinespace[2pt]
Sprendimų medis prieš Random Forest & $+$0,25 & aptikimo kokybės svoris turėtų pasiekti 0,44 (dabar 0,30); interpretuojamumo svoris turėtų pasiekti 0,01 (dabar 0,15) \\
\addlinespace[2pt]
\bottomrule
\end{tabularx}
\endgroup
